%%%%%%
%
% $Author: Sadegh Naderi $
% $Datum: 2023-10-15  $
% $Pfad: ML23-01-Keyword-Spotting-with-an-Arduino-Nano-33-BLE-Sense\Presentations\SadeghNaderi\LiteratureReview\slides\ContentLiteratureReview.tex $
% $Version: 3.0 $
% $Review Date: 2023-12-16 $
%
%
% !TeX encoding = utf8
% !TeX root = Rename
%
%%%%%%
\Mysection{Time Series Analysis: Forecasting and Control}

\STANDARD{Description}
{
	This book \cite{Box2015} is a foundational text in time series forecasting. It introduces the ARIMA modeling framework and explains methods for identifying, estimating, and validating time series models. The book provides a structured approach to modeling trend, seasonality, and autocorrelation in economic data, making it highly relevant for forecasting inflation indicators such as the Consumer Price Index (CPI).
	
	\bigskip
	
	\textbf{Limitation:} The book focuses mainly on classical statistical models and does not cover modern machine learning approaches for forecasting.
	
	\scriptsize{\textbf{Keywords:} ARIMA, Time Series Forecasting, Autocorrelation, Model Identification, Economic Forecasting}
}

\Mysection{Time Series Analysis}

\STANDARD{Description}
{
	The book \cite{Hamilton1994} presents a comprehensive econometric framework for analyzing time series data. It covers topics such as stationarity, autoregressive models, vector autoregression (VAR), cointegration, and regime-switching models. These concepts are widely used in macroeconomic forecasting and provide theoretical foundations for modeling inflation and other economic indicators.
	
	\bigskip
	
	\textbf{Limitation:} The book is mathematically intensive and may be difficult for readers without a strong background in econometrics.
	
	\scriptsize{\textbf{Keywords:} Econometrics, VAR, Cointegration, Time Series Models, Macroeconomic Forecasting}
}

\Mysection{Forecasting: Principles and Practice}

\STANDARD{Description}
{
	The textbook \cite{Hyndman2021} provides a practical introduction to forecasting techniques used in time series analysis. It explains exponential smoothing, ARIMA models, and seasonal forecasting methods in an accessible way. The book also demonstrates implementation using statistical software and highlights best practices for evaluating forecasting accuracy.
	
	\bigskip
	
	\textbf{Limitation:} While practical and accessible, the book emphasizes applied forecasting and provides less mathematical depth compared to traditional econometric texts.
	
	\scriptsize{\textbf{Keywords:} Forecasting Methods, ARIMA, Exponential Smoothing, Seasonal Forecasting, Time Series}
}

\Mysection{Time Series Analysis and Its Applications}

\STANDARD{Description}
{
	The book \cite{Shumway2017} provides a comprehensive introduction to statistical time series analysis with practical examples. It discusses spectral analysis, state-space models, ARIMA processes, and forecasting methods. The text combines theoretical foundations with applied examples, making it useful for modeling economic time series such as inflation indicators.
	
	\bigskip
	
	\textbf{Limitation:} The book uses examples primarily implemented in R, which may require adaptation for other programming environments.
	
	\scriptsize{\textbf{Keywords:} Time Series Analysis, ARIMA Models, Spectral Analysis, State-Space Models, Forecasting}
}

\Mysection{Introduction to Time Series Forecasting with Python}

\STANDARD{Description}
{
	The book \cite{Brownlee2017} provides practical guidance for implementing time series forecasting models using Python. It explains data preparation, feature engineering, and the development of forecasting models including ARIMA and machine learning methods. The book emphasizes hands-on implementation and practical workflows for forecasting tasks.
	
	\bigskip
	
	\textbf{Limitation:} The book focuses mainly on implementation and provides limited theoretical explanation of statistical forecasting models.
	
	\scriptsize{\textbf{Keywords:} Python, Time Series Forecasting, ARIMA, Data Preparation, Machine Learning}
}

\Mysection{The Elements of Statistical Learning}

\STANDARD{Description}
{
	The book \cite{Hastie2009} is a widely recognized reference for statistical learning and machine learning methods. It covers regression models, classification, neural networks, and model evaluation techniques. Although not specifically focused on time series, it provides important theoretical foundations for predictive modeling and data analysis.
	
	\bigskip
	
	\textbf{Limitation:} The book focuses primarily on general machine learning techniques and does not emphasize time series forecasting applications.
	
	\scriptsize{\textbf{Keywords:} Statistical Learning, Regression, Machine Learning, Model Evaluation, Data Mining}
}

\Mysection{Forecasting Inflation}

\STANDARD{Description}
{
	The article \cite{StockWatson1999} investigates methods for forecasting inflation using econometric time series models. The authors analyze various predictive models and evaluate their performance using macroeconomic data. The study highlights the importance of combining statistical models with economic indicators to improve inflation forecasting accuracy.
	
	\bigskip
	
	\textbf{Limitation:} The study focuses on macroeconomic forecasting models and may not directly address modern machine learning approaches for inflation prediction.
	
	\scriptsize{\textbf{Keywords:} Inflation Forecasting, Econometric Models, Macroeconomics, Predictive Modeling}
}

\Mysection{Consumer Price Index Prediction by ARIMA and ETS}

\STANDARD{Description}
{
	The article \cite{Hu2025} explores forecasting methods for the Consumer Price Index (CPI) using ARIMA and ETS models. The study evaluates the predictive performance of both approaches and demonstrates how classical time series models can effectively capture patterns in inflation data.
	
	\bigskip
	
	\textbf{Limitation:} The study focuses only on classical statistical models and does not compare them with modern machine learning forecasting approaches.
	
	\scriptsize{\textbf{Keywords:} Consumer Price Index, ARIMA, ETS, Inflation Forecasting, Time Series Models}
}

\Mysection{Forecasting Consumer Price Index: Empirical Analysis}

\STANDARD{Description}
{
	The article \cite{Mchukwa2025} presents an empirical study of forecasting Consumer Price Index values using statistical time series methods. The study analyzes historical inflation data and evaluates model accuracy using standard forecasting metrics. The results highlight the importance of model selection and proper evaluation in inflation forecasting.
	
	\bigskip
	
	\textbf{Limitation:} The analysis is based on a limited dataset and may not generalize to different economic environments.
	
	\scriptsize{\textbf{Keywords:} CPI Forecasting, Inflation Analysis, Time Series Modeling, Forecast Accuracy}
}
\Mysection{Forecasting CPI Using Machine Learning Models}

\STANDARD{Description}
{
	The article \cite{Nalici2025} compares machine learning approaches for predicting Consumer Price Index values. The study evaluates multiple algorithms and analyzes their performance in forecasting inflation trends. The results demonstrate the potential advantages of machine learning models in economic forecasting.
	
	\bigskip
	
	\textbf{Limitation:} Machine learning models may lack interpretability compared to traditional econometric forecasting methods.
	
	\scriptsize{\textbf{Keywords:} CPI Forecasting, Machine Learning, Inflation Prediction, Economic Data}
}
\Mysection{The KDD Process for Extracting Useful Knowledge from Data}

\STANDARD{Description}
{
	The article \cite{Fayyad:1996} introduces the Knowledge Discovery in Databases (KDD) process, which describes a systematic approach to extracting meaningful knowledge from large datasets. The framework includes stages such as data selection, preprocessing, transformation, data mining, and interpretation. This process provides a structured methodology for conducting data-driven research projects.
	
	\bigskip
	
	\textbf{Limitation:} The framework provides a conceptual overview but does not focus on specific algorithms or implementation techniques.
	
	\scriptsize{\textbf{Keywords:} Knowledge Discovery, Data Mining, KDD Process, Data Analysis, Pattern Discovery}
}


%-----------------------------------
